% !TEX root = ../main.tex
\begin{ouverture}[newton]{14}{Deuxième loi de Newton}{Partie 2 · Mouvement et interactions}
{Centre de masse\enspace·\enspace Référentiel galiléen\enspace·\enspace Somme des forces et accélération\enspace·\enspace Équilibre}
\objectif{Justifier qualitativement la position du centre de masse d'un système.}
\objectif{Discuter du caractère galiléen d'un référentiel pour le mouvement étudié.}
\objectif{Utiliser la deuxième loi de Newton pour déterminer l'accélération à partir des forces.}
\objectif{Utiliser la deuxième loi de Newton pour déterminer la somme des forces à partir du mouvement.}
\objectif{Exploiter la condition d'équilibre d'un système.}
\prerequis{Vecteurs vitesse et accélération, repère de Frenet (chapitre~13).}
\prerequis{Forces usuelles, principe d'inertie (classes de seconde et de première).}
\prerequis{Projection d'un vecteur sur des axes ; trigonométrie dans le triangle rectangle.}
\end{ouverture}
\essentiel
\essSection{1}{Centre de masse et référentiels}
\begin{definition}[Centre de masse]
Le \textbf{centre de masse} G d'un système est le point dont le mouvement est le plus simple : quel que soit le
mouvement de rotation du système sur lui-même, G se déplace comme un point matériel qui porterait toute la masse.
Pour un objet homogène, G est situé sur ses éléments de symétrie ; il est plus proche des parties les plus massives et
peut se trouver en dehors de la matière (anneau, boomerang).
\end{definition}
\schema{14-centre-masse}
\begin{definition}[Référentiel galiléen]
Un référentiel est \textbf{galiléen} si le principe d'inertie y est vérifié : tout système soumis à des forces qui se
compensent y est immobile ou en mouvement rectiligne uniforme. Un référentiel en translation rectiligne uniforme par
rapport à un référentiel galiléen est galiléen.
\end{definition}
\begin{propriete}[Choisir un référentiel]
Le référentiel \textbf{terrestre} peut être considéré comme galiléen pour des mouvements de courte durée devant une
journée ; le référentiel \textbf{géocentrique} pour le mouvement des satellites de la Terre ; le référentiel
\textbf{héliocentrique} pour celui des planètes. Un véhicule qui freine ou qui tourne n'est pas un référentiel
galiléen : un objet posé sans frottement s'y met en mouvement sans qu'aucune force ne l'y pousse.
\end{propriete}
\essSection{2}{La deuxième loi de Newton}
\begin{loi}[Deuxième loi de Newton]
Dans un référentiel galiléen, la somme des forces extérieures appliquées à un système de masse $m$ constante est égale
au produit de sa masse par l'accélération de son centre de masse :
\[\sum\vec{F}_\text{ext}=m\,\vec{a}_\mathrm{G}.\]
Si les forces se compensent, $\vec{a}_\mathrm{G}=\vec{0}$ : le centre de masse est immobile ou en mouvement rectiligne
uniforme (principe d'inertie). On dit que le système est \textbf{à l'équilibre} s'il est immobile.
\end{loi}
\begin{propriete}[Forces usuelles]
\par\smallskip
{\small\renewcommand{\arraystretch}{1.3}
\begin{tabularx}{\linewidth}{@{}lYl@{}}\toprule
\otTete{Force} & \otTete{Direction et sens} & \otTete{Valeur}\\\midrule
poids $\vec{P}$ & verticale, vers le bas & $P=m\,g$, $g=\qty{9,81}{N.kg^{-1}}$\\
réaction normale $\vec{N}$ & perpendiculaire au support, vers l'extérieur & déduite de la deuxième loi\\
frottement $\vec{f}$ & opposé au mouvement (ou à sa tendance) & selon la situation\\
tension $\vec{T}$ d'un fil & le long du fil, vers le point d'attache & déduite de la deuxième loi\\
force électrique & colinéaire au champ $\vec{E}$ & $F=|q|\,E$\\
gravitation & vers le centre de l'astre attracteur & $F=G\,\dfrac{m\,M}{d^2}$\\\bottomrule
\end{tabularx}}
\end{propriete}
\begin{savoirfaire}[Appliquer la deuxième loi de Newton]
1. Définir le système et le référentiel, justifier son caractère galiléen. 2. Faire le bilan des forces et les
représenter sur un schéma. 3. Écrire la deuxième loi sous forme vectorielle. 4. La projeter sur des axes bien choisis
(l'un dans le sens du mouvement, l'autre perpendiculaire). 5. En déduire l'accélération, ou une force inconnue.
\end{savoirfaire}
\begin{exemple}
Une personne de masse $m=\qty{70}{kg}$ se tient sur un pèse-personne dans un ascenseur qui démarre vers le haut avec
une accélération $a=\qty{1,5}{m.s^{-2}}$. Dans le référentiel terrestre, sur un axe vertical vers le haut :
$N-m\,g=m\,a$, d'où $N=m(g+a)=70\times(9{,}81+1{,}5)=\qty{792}{N}$. Le pèse-personne, gradué pour afficher $N/g$,
indique 80,7~kg.
\end{exemple}
\schema{14-ascenseur}
\begin{attention}
La somme des forces détermine l'\emph{accélération}, pas la vitesse : au sommet de sa trajectoire, une balle lancée
verticalement a une vitesse nulle, mais son accélération vaut toujours $\vec{g}$. Le produit $m\,\vec{a}$ n'est pas une
force : il ne figure jamais dans le bilan des forces.
\end{attention}
\essSection{3}{Des forces au mouvement}
\begin{propriete}[Prévoir un mouvement]
Connaissant les forces à chaque instant, la deuxième loi donne l'accélération ; on en déduit la vitesse puis la
position, en intégrant (chapitre~15) ou numériquement. Pour un plan incliné d'angle $\alpha$ sans frottement, par
exemple, la projection sur la ligne de plus grande pente donne $a=g\sin\alpha$, soit \qty{3,36}{m.s^{-2}} pour
$\alpha=20^\circ$.
\end{propriete}
\begin{codecote}[0.56]
\begin{pythoncol}
m, g, k, dt = 80, 9.81, 0.27, 0.1
v, t = 0, 0
while t < 30:
    a = g - k * v**2 / m    # 2e loi
    v = v + a * dt          # Euler
    t = t + dt
print(round(v, 1))
\end{pythoncol}
\begin{consolecol}
53.9
\end{consolecol}
\cote
\fontsize{9.5}{12.5}\selectfont
\capacite{Capacité numérique}\par\smallskip
Un parachutiste de 80~kg en chute libre subit son poids et une force de frottement de l'air de valeur $k\,v^2$. La
deuxième loi donne $a=g-k\,v^2/m$ ; la méthode d'Euler calcule la vitesse de proche en proche. La vitesse tend vers une
\textbf{vitesse limite}, atteinte quand le frottement compense le poids : $v_\mathrm{L}=\sqrt{m\,g/k}=\qty{53,9}{m.s^{-1}}$.
\end{codecote}
\schema{14-chute}
\begin{experience}[Vérifier la deuxième loi sur un banc à coussin d'air]
Un mobile de masse $M$ glisse sans frottement sur un banc horizontal ; il est tiré par un fil passant sur une poulie, au
bout duquel est suspendue une masse $m$. Un pointage vidéo montre que le mobile a un mouvement rectiligne uniformément
accéléré. L'accélération mesurée est conforme à la valeur $a=\frac{m\,g}{M+m}$, obtenue en appliquant la deuxième loi au
système formé du mobile et de la masse suspendue.
\end{experience}
\begin{rediger}
Toujours écrire, dans l'ordre : système, référentiel (galiléen), bilan des forces avec un schéma, loi vectorielle, puis
projections. Écrire « $\vec{P}+\vec{N}+\vec{f}=m\,\vec{a}_\mathrm{G}$ » avant de projeter, et vérifier l'homogénéité
du résultat.
\end{rediger}
\exercices

\rubrique{Découverte}
\begin{blocLivre}{ex-14-1}
\begin{exercice}[1]{Où est le centre de masse ?}
Le centre de masse G de chaque objet est indiqué sur le schéma.
\schema{14-ex01}
\begin{questions}
\item Justifier la position de G pour le disque homogène et pour l'haltère symétrique.
\item Pour l'haltère dissymétrique, les deux boules ont pour masses \qty{3,0}{kg} et \qty{1,0}{kg} ; la masse de la tige
est négligeable. Pourquoi G est-il plus près de la grosse boule ? Vérifier que $m_1\,d_1=m_2\,d_2$.
\item Que remarque-t-on pour l'anneau ? Peut-on soutenir l'anneau en équilibre en le posant par son centre de masse
sur une pointe ?
\end{questions}
\end{exercice}
\end{blocLivre}
\begin{blocLivre}{ex-14-2}
\begin{exercice}[1]{Galiléen ou non ?}
\begin{questions}[label=\otLabel{\alph*.},leftmargin=1.9em,itemsep=0pt,parsep=0.15em]
\item On étudie la chute d'une bille dans un laboratoire.
\item On étudie le mouvement d'un satellite de télécommunication autour de la Terre.
\item On étudie le mouvement de Mars autour du Soleil.
\item Dans un bus qui freine, une balle posée sur le sol roule vers l'avant.
\end{questions}
\begin{questions}
\item Pour les situations a, b et c, choisir un référentiel adapté et dire pourquoi on peut le considérer comme galiléen.
\item Dans la situation d, le référentiel du bus est-il galiléen ? Justifier à l'aide du principe d'inertie.
\item Décrire le mouvement de la balle dans le référentiel terrestre pendant le freinage.
\end{questions}
\end{exercice}
\end{blocLivre}
\begin{blocLivre}{ex-14-3}
\begin{exercice}[1]{La pierre de curling}
Une pierre de curling de masse $m=\qty{19}{kg}$ est lâchée à \qty{2,0}{m.s^{-1}} sur une piste horizontale. Pendant sa
glissade, elle subit une force de frottement constante, de valeur $f=\qty{1,9}{N}$.
\begin{questions}
\item Faire le bilan des forces appliquées à la pierre et les représenter sur un schéma.
\item Appliquer la deuxième loi de Newton et en déduire la valeur de l'accélération. Préciser son sens.
\item Quelle est la nature du mouvement ? Au bout de combien de temps la pierre s'arrête-t-elle ?
\item Quelle distance parcourt-elle ?
\end{questions}
\end{exercice}
\end{blocLivre}
\begin{blocLivre}{ex-14-4}
\begin{exercice}[1]{Le pèse-personne dans l'ascenseur}
Une personne de 60~kg est debout sur un pèse-personne dans un ascenseur qui monte. Le trajet comporte trois phases :
démarrage avec une accélération de \qty{1,2}{m.s^{-2}} vers le haut, mouvement à vitesse constante, puis arrêt avec une
accélération de \qty{1,2}{m.s^{-2}} vers le bas. Le pèse-personne affiche $N/g$, où $N$ est la force qu'il exerce sur
la personne.
\begin{questions}
\item Faire le bilan des forces appliquées à la personne.
\item Appliquer la deuxième loi de Newton sur un axe vertical orienté vers le haut.
\item Calculer la valeur affichée par le pèse-personne dans chacune des trois phases.
\item Que ressent la personne au démarrage ? à l'arrêt ?
\end{questions}
\end{exercice}
\end{blocLivre}
\begin{blocLivre}{ex-14-5}
\begin{exercice}[1]{Des systèmes en équilibre}
\begin{questions}[label=\otLabel{\alph*.},leftmargin=1.9em,itemsep=0pt,parsep=0.15em]
\item Un parachutiste équipé, de masse totale 90~kg, descend verticalement à vitesse constante.
\item Un lustre de 12~kg est suspendu au plafond par une chaîne.
\end{questions}
\begin{questions}
\item Pour chaque système, faire le bilan des forces et justifier qu'elles se compensent.
\item Calculer la force de frottement de l'air sur le parachutiste et la tension de la chaîne du lustre.
\item Le parachutiste est-il « à l'équilibre » au sens de la définition du cours ?
\end{questions}
\end{exercice}
\end{blocLivre}
\rubrique{Entraînement}
\begin{blocLivre}{ex-14-6}
\begin{exercice}[2]{Le skieur sur la pente}
Un skieur de masse $m=\qty{75}{kg}$ descend en ligne droite une piste plane inclinée de $\alpha=20^\circ$ sur
l'horizontale. Il part sans vitesse ; les frottements (neige et air) sont modélisés par une force constante de valeur
$f=\qty{120}{N}$, opposée au mouvement.
\schema{14-ex06}
\begin{questions}
\item Faire le bilan des forces et les représenter.
\item En projetant la deuxième loi de Newton sur un axe parallèle à la pente, montrer que $a=g\sin\alpha-\frac{f}{m}$.
Calculer $a$.
\item Calculer la vitesse et la distance parcourue au bout de 8,0~s.
\item En projetant sur un axe perpendiculaire à la pente, calculer la réaction normale de la neige.
\end{questions}
\end{exercice}
\end{blocLivre}
\begin{blocLivre}{ex-14-7}
\begin{exercice}[2]{Du mouvement aux forces}
Une voiture de masse $m=\qty{1300}{kg}$ passe de 0 à 100~km/h en 9,0~s sur une route horizontale rectiligne, puis
s'engage dans un rond-point de rayon \qty{25}{m} à la vitesse constante de 36~km/h.
\begin{questions}
\item En supposant l'accélération constante, calculer sa valeur pendant la première phase.
\item En déduire la valeur de la somme des forces appliquées à la voiture. Quelle force la fournit ?
\item Dans le rond-point, caractériser le vecteur accélération, puis la somme des forces.
\item Quelle force permet à la voiture de tourner ? Que se passerait-il sur une plaque de verglas ?
\end{questions}
\end{exercice}
\end{blocLivre}
\begin{blocLivre}{ex-14-8}
\begin{exercice}[2]{Le décollage d'une fusée}
Au décollage, une fusée de type Ariane 5 a une masse de \qty{7,8e5}{kg} ; ses moteurs exercent une poussée verticale de
\qty{1,3e7}{N} (valeurs arrondies). On néglige les frottements de l'air au décollage.
\begin{questions}
\item Faire le bilan des forces. Pourquoi la fusée ne peut-elle décoller que si la poussée dépasse son poids ?
\item Calculer l'accélération de la fusée au décollage.
\item En supposant cette accélération constante, calculer la vitesse et l'altitude atteintes au bout de 10~s.
\item En réalité, l'accélération augmente au cours de la montée, à poussée constante. Proposer une explication.
\end{questions}
\end{exercice}
\end{blocLivre}
\begin{blocLivre}{ex-14-9}
\begin{exercice}[2]{Une bille dans le glycérol}
Une bille d'acier de masse $m=\qty{8,4}{g}$ tombe dans du glycérol. Elle subit son poids, une poussée verticale vers le
haut de valeur constante $\Pi=\qty{0,0133}{N}$ (poussée d'Archimède, chapitre~17) et une force de frottement opposée à
sa vitesse, de valeur $f=k\,v$ avec $k=\qty{0,169}{kg.s^{-1}}$. Le programme simule sa vitesse.
\begin{python}
m, g, k, Pi = 8.4e-3, 9.81, 0.169, 0.0133      # kg, N/kg, kg/s, N
dt, v = 0.005, 0.0
for n in range(1, 61):
    a = g - Pi / m - k * v / m
    v = v + a * dt
    if n % 20 == 0:
        print(round(n * dt, 2), round(v, 3))
\end{python}
\begin{console}
0.1 0.36
0.2 0.403
0.3 0.408
\end{console}
\begin{questions}
\item \capacite{Capacité numérique}Établir l'expression de l'accélération utilisée par le programme.
\item Que vaut l'accélération au moment où la bille est lâchée sans vitesse ?
\item Interpréter les valeurs affichées. Établir l'expression de la vitesse limite et la calculer.
\item Pourquoi la bille atteint-elle sa vitesse limite beaucoup plus vite que le parachutiste de l'essentiel ?
\end{questions}
\end{exercice}
\end{blocLivre}
\begin{blocLivre}{ex-14-10}
\begin{exercice}[2]{Prendre un virage sans déraper}
Une voiture aborde un virage plat de rayon $R=\qty{50}{m}$ à vitesse constante. La force de frottement latérale exercée
par la route sur les pneus ne peut pas dépasser $\mu\,m\,g$, avec $\mu=0{,}80$ sur route sèche et $\mu=0{,}40$ sur route
mouillée.
\begin{questions}
\item Quelle force permet à la voiture de décrire le virage ? Dans quelle direction est-elle orientée ?
\item Appliquer la deuxième loi de Newton dans le repère de Frenet et exprimer la force de frottement nécessaire.
\item Montrer que la vitesse maximale vaut $v_{\max}=\sqrt{\mu\,g\,R}$ ; la calculer, en km/h, sur route sèche puis
mouillée.
\item Pourquoi cette vitesse ne dépend-elle pas de la masse du véhicule ?
\end{questions}
\end{exercice}
\end{blocLivre}
\begin{blocLivre}{ex-14-11}
\begin{exercice}[2]{Deux patineurs se repoussent}
Sur une patinoire, deux patineurs A ($m_\mathrm{A}=\qty{80}{kg}$) et B ($m_\mathrm{B}=\qty{50}{kg}$), immobiles l'un
contre l'autre, se repoussent. Juste après, B s'éloigne à \qty{2,4}{m.s^{-1}}. On néglige les frottements.
\begin{questions}
\item Faire le bilan des forces extérieures appliquées au système \{A + B\}. Que dire de son centre de masse ?
\item En déduire la vitesse de A juste après la poussée. Dans quel sens se déplace-t-il ?
\item Calculer les distances parcourues par A et par B au bout de 2,0~s, et vérifier que le centre de masse du système
n'a pas bougé.
\end{questions}
\begin{donnees}
Pour deux points de masses $m_\mathrm{A}$ et $m_\mathrm{B}$, d'abscisses $x_\mathrm{A}$ et $x_\mathrm{B}$, le centre de masse a pour abscisse
$x_\mathrm{G}=\dfrac{m_\mathrm{A}x_\mathrm{A}+m_\mathrm{B}x_\mathrm{B}}{m_\mathrm{A}+m_\mathrm{B}}$.
\end{donnees}
\end{exercice}
\end{blocLivre}
\rubrique{Objectif bac}
\begin{blocLivre}{ex-14-12}
\begin{exercice}[3]{Le saut de Felix Baumgartner}
\begin{docu}{1}{Un record}
Le 14 octobre 2012, Felix Baumgartner saute d'une capsule à 39,0~km d'altitude. Il atteint la vitesse maximale de
1357,6~km/h, franchissant le mur du son, puis ralentit avant d'ouvrir son parachute. Masse de l'ensemble sauteur et
équipement : $m=\qty{118}{kg}$.
\end{docu}
\begin{docu}{2}{Un modèle de frottement}
L'air exerce une force opposée à la vitesse, de valeur $F=\frac12\,\rho\,C_x S\,v^2$, où $\rho$ est la masse volumique de
l'air et $C_x S=\qty{0,80}{m^2}$ pour le sauteur. Masse volumique de l'air selon l'altitude :
\par\smallskip
{\small\renewcommand{\arraystretch}{1.2}
\begin{tabular}{@{}l*{5}{c}@{}}\toprule
altitude (km) & 25 & 28 & 30 & 32 & 39\\
$\rho$ (kg$\cdot$m$^{-3}$) & 0,040 & 0,025 & 0,018 & 0,013 & 0,004\\\bottomrule
\end{tabular}}
\end{docu}
\begin{donnees}
À ces altitudes, $g=\qty{9,7}{N.kg^{-1}}$.
\end{donnees}
\begin{questions}
\item Juste après le saut, la vitesse est faible et l'air très peu dense. Justifier que l'accélération vaut alors
environ $g$.
\item Dans le modèle de la chute libre, calculer la durée nécessaire pour atteindre la vitesse maximale et la hauteur
de chute correspondante.
\item \initiative Au moment où la vitesse est maximale, que vaut l'accélération ? En déduire la masse volumique de l'air
à cet instant et estimer l'altitude correspondante.
\item Expliquer pourquoi la vitesse diminue ensuite, alors que le sauteur continue de tomber.
\end{questions}
\end{exercice}
\end{blocLivre}
\begin{blocLivre}{ex-14-13}
\begin{exercice}[3]{L'accéléromètre du smartphone}
\begin{docu}{1}{Un enregistrement}
Un élève pose son smartphone sur le plancher d'un ascenseur et enregistre l'accélération verticale (axe orienté vers
le haut) pendant un trajet, l'ascenseur partant de l'arrêt.
\end{docu}
\schema{14-ex13}
\begin{questions}
\item Décrire le mouvement de l'ascenseur dans chacune des trois phases. Monte-t-il ou descend-il ?
\item Calculer la vitesse maximale atteinte et la hauteur parcourue pendant le trajet.
\item \initiative Un passager de 65~kg se tient debout. Calculer la force exercée par le plancher sur lui dans chaque
phase.
\item Le smartphone mesure en réalité $a_z+g$ lorsqu'il est immobile. Pourquoi, selon vous, son logiciel retranche-t-il
$g$ pour afficher l'enregistrement ci-dessus ?
\end{questions}
\end{exercice}
\end{blocLivre}
\begin{blocLivre}{ex-14-14}
\begin{exercice}[3]{Le téléski}
\begin{docu}{1}{Remonter la pente}
Un skieur de masse $m=\qty{70}{kg}$ est tiré par la perche d'un téléski sur une piste rectiligne inclinée de
$\alpha=15^\circ$. La perche fait un angle $\beta=25^\circ$ avec la piste. Les frottements de la neige sont modélisés par
une force constante de valeur $f=\qty{60}{N}$, parallèle à la piste.
\end{docu}
\schema{14-ex14}
\begin{questions}
\item Faire le bilan des forces appliquées au skieur et les représenter.
\item Le skieur monte à vitesse constante. Que peut-on dire de la somme des forces ? Justifier.
\item \initiative Calculer la force exercée par la perche, puis la réaction normale de la neige.
\item Au départ, le skieur passe de 0 à \qty{2,5}{m.s^{-1}} en 2,0~s. Quelle force la perche exerce-t-elle alors ?
\end{questions}
\end{exercice}
\end{blocLivre}
\begin{blocLivre}{ex-14-15}
\begin{exercice}[3]{Le lancer du marteau}
\begin{docu}{1}{Un engin de 7,26 kg}
Le marteau est une boule métallique reliée à une poignée par un câble d'acier ; sa masse totale vaut \qty{7,26}{kg},
dont environ \qty{7,0}{kg} pour la boule. Juste avant le lâcher, le centre de masse du marteau décrit, à vitesse
constante, un cercle de rayon $R=\qty{1,9}{m}$ ; sa vitesse atteint \qty{27}{m.s^{-1}}.
\end{docu}
\begin{docu}{2}{Après le lâcher}
Une fois lâché, le marteau tourne sur lui-même pendant son vol ; les frottements de l'air sont négligeables.
\end{docu}
\begin{questions}
\item Où se trouve le centre de masse du marteau ? Justifier qualitativement.
\item Avant le lâcher, caractériser l'accélération du centre de masse, puis la force exercée par le câble.
\item \initiative Comparer cette force au poids d'un lanceur de 110~kg. Pourquoi le lanceur se penche-t-il vers
l'arrière en tournant ?
\item Après le lâcher, quelle est la seule force appliquée au marteau ? Que peut-on dire de l'accélération de son centre
de masse, et de sa trajectoire ?
\end{questions}
\begin{donnees}
$g=\qty{9,81}{N.kg^{-1}}$.
\end{donnees}
\end{exercice}
\end{blocLivre}
